\newpage
\setcounter{dang}{0}
\section{QUY TẮC TÍNH ĐẠO HÀM}
\subsection{KIẾN THỨC CẦN NHỚ}

\subsubsection{Quy tắc tính đạo hàm của tổng hiệu, tích thương}
Giả sử $u=u(x)$, $v=v(x)$ là các hàm số có đạo hàm tại điểm $x$ thuộc khoảng xác định. Ta có các quy tắc sau:
\begin{tcolorbox}[colframe=cyan,colback=red!1!white,boxrule=0.3mm]
	\begin{itemize}
		\begin{multicols}{2}
			\item [\ding{172}] $(u+v)'=u'+v'$
			\item [\ding{173}] $(u-v)'=u'-v'$
			\item [\ding{174}] $(u\cdot v)'=u'\cdot v+v' \cdot u$
			\item [\ding{175}] $\left( \dfrac{u}{v}\right)'=\dfrac{u' \cdot v-v' \cdot u}{v^2}$,  với $v \ne 0$.
			\item [\ding{176}] $(k \cdot u)'=k \cdot u'$, với $k$ là hằng số.
			\item [\ding{177}] $\left( \dfrac{1}{v}\right)'=-\dfrac{v'}{v^2}$, với $v \ne 0$.
		\end{multicols}
	\end{itemize}
\end{tcolorbox}

\subsubsection{Đạo hàm của hàm hợp}
Cho hàm số $y=f(u)$, với $u=u(x)$. Ta có công thức sau:
\boxmini{$[f(u)]'=f'(u) \cdot u'_x$}
% \subsubsection{Bảng đạo hàm của một số hàm số sơ cấp cơ bản và hàm hợp}
% \begin{center}
% \begin{tikzpicture}[xscale=8,yscale=1]
% 	\begin{scope}[shift={(-.5,.5)}]
% 		\fill[pink!5] (0,-1) rectangle (2,-12);
% 		\fill[cyan!25] (0,0) rectangle (2,-1);
% 		\draw [line width=0.5pt](0,0) grid (2,-12)
% 		(1,-1)--(2,-1) (1,-2)--(2,-2)
% 		(1,-3)--(2,-3) (0,-4)--(2,-4)
% 		(1,-5)--(2,-5) (0,-6)--(2,-6)
% 		(1,-7)--(2,-7) (0,-8)--(2,-8)
% 		(1,-9)--(2,-9) (1,-10)--(2,-10)
% 		(1,-11)--(2,-11) (0.15,0)--(0.15,-12)
% 		;
% 	\end{scope}
% 	% \path
% 	% (-0.42,0) node{\text{\textbf{STT}}}
% 	% (-0.42,-1) node{\circEX{1}} 
% 	% (-0.42,-2) node{\circEX{2}}
% 	% (-0.42,-3) node{\circEX{3}}
% 	% (-0.42,-4) node{\circEX{4}}
% 	% (-0.42,-5) node{\circEX{5}}
% 	% (-0.42,-6) node{\circEX{6}}
% 	% (-0.42,-7) node{\circEX{7}}
% 	% (-0.42,-8) node{\circEX{8}}
% 	% (-0.42,-9) node{\circEX{9}}
% 	% (-0.42,-10) node{\circEX{10}}
% 	% (-0.42,-11) node{\circEX{11}} 
% 	% ;
% 	\path
% 	(0.08,0) node{\text{\textbf{Đạo hàm của hàm sơ cấp cơ bản}}}  
% 	(-0.2,-1) node[right]{$\left(x^n\right)'=n \cdot x^{n-1}$}
% 	(-0.2,-2) node[right]{$\left({\dfrac{1}{x}}\right)'=-\dfrac{1}{x^2}$}    
% 	(-0.2,-3) node[right]{$\left(\sqrt{x}\right)'=\dfrac{1}{2\sqrt{x}}$}
% 	(-0.2,-4) node[right]{$\left({\sin x}\right)'=\cos x$}    
% 	(-0.2,-5) node[right]{$\left({\cos x}\right)'=-\sin x$}
% 	(-0.2,-6) node[right]{$\left({\tan x}\right)'=\dfrac{1}{{\cos}^2x}$}    
% 	(-0.2,-7) node[right]{$\left({\cot x}\right)'=-\dfrac{1}{{\sin}^2x}$}
% 	(-0.2,-8) node[right]{$\left({e^x}\right)'=e^x$}    
% 	(-0.2,-9) node[right]{$\left({a^x}\right)'=a^x.\ln a$}
% 	(-0.2,-10) node[right]{$\left({\ln x}\right)'=\dfrac{1}{x}$}
% 	(-0.2,-11) node[right]{$\left({\log_a x}\right)'=\dfrac{1}{x \cdot \ln a}$}
% 	(1,0) node{\text{\textbf{Đạo hàm của hàm hợp, với u = u(x)}}}    
% 	(0.8,-1) node[right]{$\left(u^n\right)'=n \cdot u^{n-1} \cdot u'$}
% 	(0.8,-2) node[right]{$\left({\dfrac{1}{u}}\right)'=-\dfrac{u'}{u^2}$}    							
% 	(0.8,-3) node[right]{$\left(\sqrt{u}\right)'=\dfrac{u'}{2\sqrt{u}}$}
% 	(0.8,-4) node[right]{$\left({\sin u}\right)'=u' \cdot \cos u$}    
% 	(0.8,-5) node[right]{$\left({\cos u}\right)'=-u' \cdot \sin u$}
% 	(0.8,-6) node[right]{$\left({\tan u}\right)'=\dfrac{u'}{{\cos}^2u}$}    
% 	(0.8,-7) node[right]{$\left({\cot u}\right)'=-\dfrac{u'}{{\sin}^2u}$}
% 	(0.8,-8) node[right]{$\left({e^u}\right)'=u' \cdot e^u$}    
% 	(0.8,-9) node[right]{$\left({a^u}\right)'=u' \cdot a^u.\ln a$}
% 	(0.8,-10) node[right]{$\left({\ln u}\right)'=\dfrac{u'}{u}$}
% 	(0.8,-11) node[right]{$\left({\log_a u}\right)'=\dfrac{u'}{u \cdot \ln a}$}
% 	;
% \end{tikzpicture}
% \end{center}
\subsubsection*{Bảng đạo hàm}
\begin{center}
	\renewcommand{\arraystretch}{2}
	\begin{longtable}{|C{0.45\textwidth}|C{0.45\textwidth}|}
	\hline
	\centerline{\bf Đạo hàm của hàm số sơ cấp cơ bản}& \centerline{\bf Đạo hàm của hàm hợp}\\
	\centerline{\bf thường gặp} & \centerline{\bf(ở đây $\mathbf{u=u(x)}$)}\\
	\hline
	$\left(x^n\right)'=n\cdot x^{n -1}$ & $\left(u^n\right)'=n\cdot u^{n-1}\cdot u'$\\
	\hline
	$\left(\dfrac{1}{x}\right)'=-\dfrac{1}{x^2}$& $\left(\dfrac{1}{u}\right)'=-\dfrac{u'}{u^2}$\\
	\hline
	$\left(\sqrt{x}\right)'=\dfrac{1}{2\sqrt{x}}$&
	$\left(\sqrt{u}\right)'=\dfrac{u'}{2\sqrt{u}}$\\
	\hline
	$\left(\sin x\right)'=\cos x$& $\left(\sin u\right)'=u'\cdot \cos u$\\
	\hline
	$\left(\cos x\right)'=-\sin x$& $\left(\cos u\right)'=-u'\cdot\sin u $ \\
	\hline
	$\left(\tan x\right)'=\dfrac{1}{\cos^2 x}$&$\left(\tan u\right)'=\dfrac{u'}{\cos^2 u}$ \\
	\hline
	$\left(\cot x\right)'=-\dfrac{1}{\sin^2 x}$&$\left(\cot u\right)'=-\dfrac{u'}{\sin^2 u}$ \\
	\hline
	$\left(\mathrm{e}^x\right)'=\mathrm{e}^{x}$& $\left(\mathrm{e}^{u}\right)'=u'\cdot \mathrm{e}^{u}$ \\
	\hline
	$\left(a^{x}\right)'=a^{x}\ln a$&$\left(a^{u}\right)=u'\cdot a^{u}\ln a$ \\
	\hline
	$\left(\ln x\right)'=\dfrac{1}{x}$& $\left(\ln u\right)'=\dfrac{u'}{u}$ \\
	\hline
	$\left(\log_{a}x\right)' =\dfrac{1}{x\ln a}$& $\left(\log_{a} u\right)'=\dfrac{u'}{u\ln a}$\\
	\hline
	\end{longtable}
\end{center}
\subsection{PHÂN LOẠI, PHƯƠNG PHÁP GIẢI TOÁN}
\begin{dang}{Tính đạo hàm của hàm đa thức}
	% \begin{itemize}
	% 	\item [\iconCH] Công thức sử dụng: Với $n \in \mathbb{N}^*$, ta có:
	% 	\begin{listEX}[2]
	% 		\item [\faCaretRight] \fbox{$\big(x^n\big)'=n x^{n-1}$}
	% 		\item [\faLongArrowRight] Hàm hợp: \fbox{$\big(u^n\big)'=nu^{n-1}\cdot u'$}
	% 	\end{listEX}
	% \end{itemize}
\end{dang}

\begin{vd}
	Tính đạo hàm của các hàm số sau:
	\begin{tasks}(1)
		\task $y=x^5$;
		\task $y=2x^4$;
		\task $y=\dfrac{1}{3}x^3+\sqrt{3}$;
		\task $y=x^3+x^2-6$;
		\task $y=2x^3-3x^2+1$;
		\task $y=x^5(3-2x^4)$.
	\end{tasks}
\end{vd} \dongcham{7}

\begin{vd}
	Tính đạo hàm của các hàm số sau:
	\begin{tasks}(1)
		\task $y=(2x-3)^2$;
		\task $y=(x^2+x)^8$;
		\task $y=(x^2-5x^4)^6$.
	\end{tasks}
\end{vd} \dongcham{7}

\begin{vd}
	Giải phương trình $f'(x)=0$, với
	\begin{tasks}(1)
		\task $f(x)=x^3+3x^2-9x+5$;
		\task $f(x)=x^3+3x^2$;
		\task $f(x)=-x^4+6x^2-3$;
		\task $f(x)=\dfrac{x^4}{4}-x^3+3x^2-4x+4$.
	\end{tasks}
\loigiai{
	\begin{enumerate}[a)]
		\item
		\item
		\item
		\item
\end{enumerate}}
\end{vd} \dongcham{7}

\begin{vd}
	Cho hàm số $f(x)=x^{3}+3x^{2}-9x+1$. Giải bất phương trình $f'(x)>0$.
	\loigiai{
		Ta có 
		\begin{itemize}
			\item Tập xác định $\mathbb{R}$.
			\item $f'(x)=3x^{2}+6x-9.$
			\item $f'(x)>0  \Leftrightarrow  3x^{2}+6x-9>0
			\Leftrightarrow  \hoac{&x<-3\\&x>1.}$
		\end{itemize}
		Vậy phương trình đã cho có tập nghiệm là $S=(-\infty;-3)\cup(1;+\infty).$		
	}
\end{vd} \dongcham{4}

\begin{vd}
	Cho hàm số $ f(x) = \dfrac{x^3}{3} + \dfrac{x^2}{x} +x $. Giải bất phương trình $ f'(x) \le 0 $.
	\loigiai{
		Ta có: $y' = x^2 + x + 1 = \left(x+\dfrac{1}{2}\right)^2 + \dfrac{3}{4} > 0 \quad \forall x .$ \\
		Vậy tập nghiệm của bất phương trình $ f'(x) \le 0 $ là: $ \mathcal{S} = \varnothing$.}
\end{vd} \dongcham{4}

\begin{dang}{Tính đạo hàm của hàm chứa căn thức}
	% \begin{itemize}
	% 	\item [\iconCH] Công thức sử dụng: Với $x>0$, ta có:
	% 	\begin{listEX}[2]
	% 		\item [\faCaretRight] \fbox{$\big(\sqrt{x}\big)'=\dfrac{1}{2\sqrt{x}}$}
	% 		\item [\faLongArrowRight] Hàm hợp: \fbox{$\big(\sqrt{u}\big)'=\dfrac{u'}{2\sqrt{u}}$}
	% 	\end{listEX}
	% \end{itemize}
\end{dang}

\begin{vd}
	Tính đạo hàm của các hàm số sau:
	\begin{tasks}(1)
		\task $y=x^2+2\sqrt{x}$;
		\task $y=\dfrac{x^3}{3}+\dfrac{1}{\sqrt{x}}+\sqrt{3}$;
		\task $y=x\sqrt{x}$.
	\end{tasks}
\end{vd} \dongcham{7}

\begin{vd}
	Tính đạo hàm của các hàm số sau:
	\begin{tasks}(1)
		\task $y=\sqrt{x^2-2x+7}$;
		\task $y=\sqrt{4-x^2}$;
		\task $y=\sqrt{2x^2(2-x^2)}$.
	\end{tasks}
\end{vd} \dongcham{10}

\begin{vd}
	Giải phương trình $f'(x)=0$ với
	\begin{tasks}(1)
		\task $f(x)=\sqrt{x^2-4x+7}$;
		\task $f(x)=\sqrt{x^2-5x+6}$;
		\task $f(x)=x+\sqrt{x^2-5x+6}$
	\end{tasks}
\end{vd} \dongcham{12}

\begin{dang}{Tính đạo hàm của hàm lượng giác}
	% \begin{itemize}
	% 	\item [\iconCH] Công thức sử dụng:
	% 	\begin{listEX}[2]
	% 		\item [\faCaretRight] \fbox{$\big(\sin x\big)'=\cos x$.}
	% 		\item [\faLongArrowRight] Hàm hợp: \fbox{$\big(\sin u\big)'=u' \cdot \cos u$.}
	% 	\end{listEX}
	% 	\begin{listEX}[2]
	% 		\item [\faCaretRight] \fbox{$\big(\cos x\big)'=-\sin x$.}
	% 		\item [\faLongArrowRight] Hàm hợp: \fbox{$\big(\cos u\big)'=-u' \cdot \sin u$.}
	% 	\end{listEX}
	% 	\begin{listEX}[2]
	% 		\item [\faCaretRight] \fbox{$\left( \tan x\right)'=\dfrac{1}{\cos^2x}$.}
	% 		\item [\faLongArrowRight] Hàm hợp: \fbox{$\left( \tan u\right)'=\dfrac{u'}{\cos^2u}$.}
	% 	\end{listEX}
	% 	\begin{listEX}[2]
	% 		\item [\faCaretRight] \fbox{$\left( \cot x\right)'=-\dfrac{1}{\sin^2x}$}
	% 		\item [\faLongArrowRight] Hàm hợp: \fbox{$\left( \cot u\right)'=-\dfrac{u'}{\sin^2u}$.}
	% 	\end{listEX}
	% \end{itemize}
\end{dang}

\begin{vd}
	Tính đạo hàm của các hàm số sau:
	\begin{tasks}(1)
		\task $y=\sin x +x$;
		\task $y=\sin x +3\cos x$;
		\task $y=\tan x -\sqrt{3} \cot x$;
		\task $y=\sin (2x)$;
		\task $y=\sin^2 x$;
		\task $y=\cos^2 4x$;
	\end{tasks}
\loigiai{
	\begin{enumerate}[a)]
		\item $y'=\cos x +1$.
		\item $y'=\cos x -3\sin x$.
		\item $y'=\dfrac{1}{\cos^2 x} +\dfrac{\sqrt{3}}{\sin^2 x}$.
		\item $y'=2\cos 2x$.
		\item $y'=2\sin x \cos x=\sin 2x$.
		\item $y'=2\cos 4x \sin 4x \cdot 4 = 4\sin 8x$.
	\end{enumerate}}
\end{vd} \dongcham{10}

\begin{dang}{Tính đạo hàm của hàm số mũ, hàm số lôgarit}
	% \begin{itemize}
	% 	\item [\iconCH] Công thức sử dụng:
	% \begin{listEX}[2]
	% 	\item [\faCaretRight] \fbox{$\big(\mathrm{e}^x\big)'=\mathrm{e}^ x$.}
	% 	\item [\faLongArrowRight] Hàm hợp: \fbox{$\big(\mathrm{e}^u\big)'=u' \cdot \mathrm{e}^ u$.}
	% \end{listEX}
	% \begin{listEX}[2]
	% 	\item [\faCaretRight] \fbox{$\big(a^x\big)'=a^ x \cdot \ln a$.}
	% 	\item [\faLongArrowRight] Hàm hợp: \fbox{$\big(a^u\big)'=u' \cdot a^ u \cdot \ln a$.}
	% \end{listEX}
	% \begin{listEX}[2]
	% 	\item [\faCaretRight] \fbox{$\left({\ln x}\right)'=\dfrac{1}{x}$.}
	% 	\item [\faLongArrowRight] Hàm hợp: \fbox{$\left({\ln u}\right)'=\dfrac{u'}{u}$.}
	% \end{listEX}
	% \begin{listEX}[2]
	% 	\item [\faCaretRight] \fbox{$\left({\log_a x}\right)'=\dfrac{1}{x \cdot \ln a}$.}
	% 	\item [\faLongArrowRight] Hàm hợp: \fbox{$\left({\log_a u}\right)'=\dfrac{u'}{u \cdot \ln a}$.}
	% \end{listEX}
	% \end{itemize}
\end{dang}

\begin{vd}
	Tính đạo hàm của các hàm số sau:
	\begin{tasks}(1)
		\task $y=10^x$
		\task $y=\log_5x$ 
		\task $y=\ln (1+x^2)$ 
		\task $y=\mathrm{e}^{2x+1}$ 
		\task $y=\log_2\left(x+\mathrm{e}^x\right)$.
		\task $y=\log_8\left(x^2-3x-4\right)$
	\end{tasks}
\loigiai{
	\begin{enumerate}[a)]
		\item $\left( {10}^{x} \right)'={10}^{x}\ln10$.
		\item $y'=\dfrac{1}{x\ln5}$
		\item $ y'=\dfrac{(1+x^2)'}{1+x^2}=\dfrac{2x}{1+x^2}$.
		\item $y'=2\cdot \mathrm{e}^{2x+1}$.
		\item $y'=\dfrac{\left(x+\mathrm{e}^x\right)'}{\left(x+\mathrm{e}^x\right)\ln 2}=\dfrac{1+\mathrm{e}^x}{\left(x+\mathrm{e}^x\right)\ln 2}$.
		\item $y'=\left[\log_8\left(x^2-3x-4\right)\right]'=\dfrac{\left(x^2-3x-4\right)'}{\left(x^2-3x-4\right)\ln 8}=\dfrac{2x-3}{\left(x^2-3x-4\right)\ln 8}$.
	\end{enumerate}
	
}
\end{vd} \dongcham{16}

\begin{vd}
	Cho hàm số $f(x)=\ln \left(x+\sqrt{x^2+1}\right)$. Tính $f'(1)$.
	\loigiai{
		Ta có: $f\left(x\right)=\ln \left(x+\sqrt{x^2+1}\right) \Rightarrow f'\left(x\right)=\dfrac{\left(x+\sqrt{x^2+1}\right)'}{x+\sqrt{x^2+1}}=\dfrac{1+\dfrac{x}{\sqrt{x^2+1}}}{x+\sqrt{x^2+1}}=\dfrac{1}{\sqrt{x^2+1}}$.\\
		Vậy $f'\left(1\right)=\dfrac{1}{\sqrt{2}}$.
	}
\end{vd} \dongcham{8}

\begin{dang}{Tính đạo hàm dạng tích hoặc thương}
		\indamm{Công thức giải nhanh:}
		% \begin{boxdl}
			\begin{itemize}
				\item [\ding{172}]  $\left(\dfrac{ax+ b}{cx + d}\right)' = \dfrac{ad - bc}{\left(cx + d\right)^2}$.
				\item [\ding{173}]  $\left(\dfrac{ax^2+bx+c}{dx+e}\right)' = \dfrac{adx^2+2aex+be-cd}{(dx+e)^2}$
				\item [\ding{174}]  $\left(\dfrac{a_1 x^2 + b_1 x + c_1}{a_2 x^2 + b_2 x + c_2}\right)' = \dfrac{\begin{vmatrix}
						a_1&b_1\\
						a_2&b_2
					\end{vmatrix}x^2 + 2 \cdot \begin{vmatrix}
						a_1&c_1\\
						a_2&c_2
					\end{vmatrix}x + \begin{vmatrix}
						b_1&c_1\\
						b_2&c_2
				\end{vmatrix}}{\left(a_2 x^2 + b_2 x + c_2\right)^2}$
				\begin{itemize}
					\item [$\bullet$] Định thức trên tử số, ta nhân theo chiều "dấu huyền" trừ cho "dấu sắc".
					\item [$\bullet$] Nghĩa là: $\begin{vmatrix}
						a_1&b_1\\
						a_2&b_2
					\end{vmatrix}=a_1b_2-a_2b_1.$
				\end{itemize}
			\end{itemize}
\end{dang}

\begin{vd}
	Tính đạo hàm của các hàm số sau:
	\begin{tasks}(1)
		\task $y = \dfrac{2x + 1}{x + 1}$
		\task $y = \dfrac{2x - 1}{4x - 3}$
		\task $y = \dfrac{5x}{1 - 4x}.$
		\task $y = \dfrac{x^2 - 3x + 7}{2x - 5}$
		\task $y = \dfrac{2x - 5}{x^2 + x + 2}$
		\task $y = \dfrac{1 + x - x^2}{1 - x + x^2}.$
	\end{tasks}
	\loigiai{\begin{enumerate}[a)]
			\item 
			\begin{itemize}
				\item Sử dụng công thức $\left(\dfrac{u}{v}\right)' = \dfrac{u' \cdot v - v' \cdot u}{v^2}$, ta được:\\
				$y' = \dfrac{\left(2x + 1\right)' \left(x+1\right) - \left(x+1\right)' \left(2x+1\right)}{\left(x+1\right)^2} = \dfrac{2\left(x+1\right)-1\left(2x+1\right)}{\left(x+1\right)^2} = \dfrac{1}{\left(x+1\right)^2}.$
				\item \textbf{Cách 2:} Sử dụng công thức nhanh (rèn luyện cho lớp $12$ trong trắc nghiệm):\\
				$y ' = \dfrac{2\cdot 1 - 1 \cdot 1}{\left(x+1\right)^2} = \dfrac{1}{\left(x+1\right)^2}$
			\end{itemize}
			\item $y'=\dfrac{-2}{\left(2x-1\right)^2}$.
			\item $y'=\dfrac{5}{\left(4x-1\right)^2}$.
			\item $y'=\dfrac{2x^2 - 10x + 1}{\left(2x-5\right)^2}.$
			\item $y'=-\dfrac{2x^2 - 10x - 9}{\left(x^2 + x + 2\right)^2}$
			\item $y'=-\dfrac{2\left(2x-1\right)}{\left(x^2-x+1\right)^2}.$
	\end{enumerate}}
\end{vd} \dongcham{18}

\begin{vd}
	Tính đạo hàm của các hàm số sau:
	\begin{tasks}(1)
		\task $y=x^2\sqrt{x}$.
		\task $y=(x-1)\sqrt{3x+2}$.
		\task $y=x\sqrt{1-x^2}$
	\end{tasks}
	\loigiai{
		\begin{enumerate}[a)]
			\item $y'=(x^2)'\cdot \sqrt{x}+x^2\cdot (\sqrt{x})'=2x\sqrt{x}+\dfrac{1}{2}x\sqrt{x}=\dfrac{5}{2}x\sqrt{x}$.
			\item $y'=(x-1)'\sqrt{3x+2}+(x-1)\left(\sqrt{3x+2}\right)'=\sqrt{3x+2}+(x-1)\cdot \dfrac{3}{2\sqrt{3x+2}}=\dfrac{2(3x+2)+3(x-1)}{2\sqrt{3x+2}}=\dfrac{9x+1}{2\sqrt{3x+2}}$.
			\item $y'=\sqrt{1-x^2}+x\cdot\dfrac{(1-x^2)'}{2\sqrt{1-x^2}}=\sqrt{1-x^2}-\dfrac{x^2}{\sqrt{1-x^2}}=\dfrac{1-2x^2}{\sqrt{1-x^2}}$.	
			\end{enumerate}}
\end{vd} \dongcham{15}

\begin{vd}
	Tính đạo hàm của các hàm số sau:
	\begin{tasks}(1)
		\task $y=x\cdot \cos x$.
		\task $y=(2-x^2)\cdot \sin 3x$.
		\task $y=\dfrac{x}{\cos x}$.
		\task $y=\dfrac{\sin x-\cos x}{\sin x+\cos x}$.
	\end{tasks}
	\loigiai{
		\begin{enumerate}[a)]
			\item Ta có
			\begin{eqnarray*}
				y'&=& (x)'\cdot \cos x + x\cdot (\cos x)'\\
				&=& \cos x - x\cdot \sin x.
			\end{eqnarray*}
			\item 
			\begin{eqnarray*}
				y'&=& (2-x^2)'\cdot \sin 3x + (2-x^2)\cdot (\sin 3x)'\\
				&=& -2x\cdot \sin 3x + (2-x^2)\cdot 3\cos 3x.
			\end{eqnarray*}
			\item 
			\begin{eqnarray*}
				y'&=& \dfrac{(x)'\cdot \cos x - x\cdot (\cos x)'}{\cos^2 x}\\
				&=& \dfrac{\cos x + x\cdot \sin x}{\cos^2 x}.
			\end{eqnarray*}
			\item Ta có
			\begin{eqnarray*}
				y'&=& \dfrac{(\sin x-\cos x)'\cdot (\sin x+\cos x)-(\sin x-\cos x)\cdot (\sin x+\cos x)'}{(\sin x+\cos x)^2}\\
				&=& \dfrac{(\cos x+\sin x)\cdot (\sin x+\cos x)-(\sin x-\cos x)\cdot (\cos x-\sin x)}{(\sin x+\cos x)^2}\\
				&=& \dfrac{(\cos x+\sin x)^2+(\sin x-\cos x)^2}{(\sin x+\cos x)^2}\\
				&=& \dfrac{(1+2\sin x\cdot \cos x)+(1-2\sin x\cdot \cos x)}{(\sin x+\cos x)^2}\\
				&=& \dfrac{2}{(\sin x+\cos x)^2}.
			\end{eqnarray*}
		\end{enumerate}}
\end{vd} \dongcham{20}

\begin{dang}{Viết phương trình tiếp tuyến}
	% \subsubsection{Đề bài cho trước $(x_0;y_0)$:}
	% \begin{itemize}
	% 	\item [\ding{172}] Tính $f'(x_0)$.
	% 	\item [\ding{173}] Thay vào công thức \fbox{$y=f'(x_0)(x-x_0)+y_0$}, thu gọn kết quả về dạng $y=Ax+B$.
	% \end{itemize}
	% \subsubsection{Đề bài chưa cho đầy đủ $(x_0;y_0)$, ta thường gặp các loại sau:}
	% \begin{itemize}
	% 	\item [\ding{172}] Cho biết trước $x_0$ hoặc $y_0$. Ta chỉ việc thay giá trị đó vào hàm số $y=f(x)$, sẽ tính được đại lượng còn lại.
	% 	\item [\ding{173}] Cho trước 1 điều kiện giải. Ta chỉ việc giải điều kiện đó, tìm $x_0$.
	% 	\begin{itemize}
	% 		\item [$\bullet$] Nếu đề bài cho biết hệ số góc của tiếp tuyến $k=a$, ta giải $f'(x)=a$ tìm nghiệm $x_0$.
	% 		Đề bài thường cho hệ số góc tiếp tuyến dưới các dạng sau:
	% 		\begin{tcolorbox}[colframe=cyan,colback=red!1!white,boxrule=0.1mm]
	% 			\begin{itemize}
	% 				\item Tiếp tuyến $d\parallel\Delta\colon y=ax+b\Rightarrow k=a$
	% 				\item Tiếp tuyến $d\perp\Delta\colon y=ax+b\Rightarrow k=-\dfrac{1}{a}$
	% 			\end{itemize}
	% 		\end{tcolorbox}
			
	% 		\item [$\bullet$] Nếu đề bài cho tiếp điểm là giao điểm của hai đồ thị $y=f(x)$ và $y=g(x)$, ta giải $f(x)=g(x)$ để tìm nghiệm $x_0$.
	% 	\end{itemize}
	% \end{itemize}
\end{dang}

\begin{vd}
	Cho đường cong $(C):y=f(x)=\dfrac{x^2}{2}-4x+1$.
	\begin{tasks}(1)
		\task Viết phương trình tiếp tuyến của $(C)$ tại điểm có hoành độ $x_0=-2$.
		\task Viết phương trình tiếp tuyến của $(C)$ tại điểm có hoành độ $x_0=-3$.
		\task Viết phương trình tiếp tuyến của $(C)$, biết tiếp tuyến có hệ số góc $k=1$.
	\end{tasks}
	\loigiai{
		\begin{enumerate}[a)]
			\item Ta có $f'(x)=x-4$. Với $x_0=-2\Rightarrow y_0=11$.\\
			Do đó, tiếp tuyến cần tìm có phương trình: $y=f'(-2)(x+2)+11=-6x-1$.
			\item Gọi $(x_0;y_0)$ là tiếp điểm. Ta có $f'(x_0)=1\Leftrightarrow x_0-4=1\Leftrightarrow x_0=5\Rightarrow y_0=-\dfrac{13}{2}$.\\
			Vậy, tiếp tuyến có phương trình là $y=1(x-5)-\dfrac{13}{2}=x-\dfrac{23}{2}$.
	\end{enumerate}}
\end{vd} \dongcham{16}

\begin{vd}
	Viết phương trình tiếp tuyến của đường cong $(C):y=f(x)=x(x^2+x-1)+1$ tại điểm có tung độ bằng $-1$.
	\loigiai{
		Đạo hàm $y'=3x^2+2x-1$. Gọi $(x_0;y_0)$ là tiếp điểm, ta có $$y_0=-1\Leftrightarrow x_0(x_0^2+x_0-1)+1=-1\Leftrightarrow (x_0+2)(x_0^2-x_0+1)=0\Leftrightarrow x_0=-2.$$
		Tính được $f'(-2)=7$, ta có phương trình tiếp tuyến cần tim: $y=7x+13$.}
\end{vd} \dongcham{10}

\begin{vd}
	Cho hàm số $y=f(x)=x^3-3x^2+2$ có đồ thị $(C)$. Viết phương trình tiếp tuyến của $(C)$ biết tiếp tuyến song song với đường thẳng $\Delta:3x+y=2$.
	\loigiai{Gọi $(x_0;y_0)$ là tiếp điểm của tiếp tuyến cần tìm.\\
		Vì tiếp tuyến song song với đường thẳng $\Delta:y=-3x+2$ nên ta có $$f'(x_0)=-3\Leftrightarrow 3x_0^2-6x_0+3=0\Leftrightarrow x_0=1\Rightarrow y_0=0.$$
		Do vậy, tiếp tuyến có phương trình: $y=-3(x-1)+0=-3x+3$.}
\end{vd} \dongcham{12}

\begin{vd}
	Viết phương trình tiếp tuyến của đồ thị $(C):y=\dfrac{x-1}{x+2}$ biết tiếp tuyến vuông góc với đường thẳng $\Delta:3x+y-2=0$.
	\loigiai{
		Tập xác định: $\mathscr{D}=\mathbb{R}\setminus\{-2\}$. Đạo hàm $y'=\dfrac{3}{(x+2)^2}$. Viết lại phương trình đường thẳng $\Delta:y=-3x+2$. Gọi $(x_0;y_0)$ là tiếp điểm của tiếp tuyến cần tìm. Tiếp tuyến vuông góc với đường thẳng $\Delta$ nên $$f'(x_0)=\dfrac{1}{3}\Leftrightarrow \dfrac{3}{(x_0+2)^2}=\dfrac{1}{3}\Leftrightarrow \hoac{&x_0=1\\&x_0=-5.}$$
		Từ đó tìm được hai tiếp tuyến thỏa mãn yêu cầu bài toán là: $y=\dfrac{x-1}{3}$ và $y=\dfrac{x+11}{3}$.}
\end{vd} \dongcham{12}

\begin{vd}
	Viết phương trình tiếp tuyến $\Delta$ của đồ thị $(C):y=x^2-3x+2$ biết tiếp tuyến $\Delta$ cắt trục tung tại điểm có tung độ bằng $-2$.
	\loigiai{
		Gọi $(x_0;y_0)$ là tiếp điểm của tiếp tuyến $\Delta$. Vì $\Delta$ cắt trục tung tại điểm có tung độ bằng $-2$ nên phương trình tiếp tuyến $\Delta$ có dạng $y=kx-2$.\\
		Ta có $f'(x)=2x-3$. Do đó, $$f'(x_0)=k\Leftrightarrow 2x_0-3=k.$$
		Mặt khác, vì $(x_0;y_0)$ thuộc đồ thị $(C)$ nên $y_0=x_0^2-3x_0+2$. Thay vào phương trình tiếp tuyến, ta được $$x_0^2-3x_0+2=(2x_0-3)x_0-2\Leftrightarrow x_0^2-3x_0+2=2x_0^2-3x_0-2\Leftrightarrow x_0^2=4\Leftrightarrow x_0=\pm 2.$$
		Với $x_0=2$, ta có $k=1$ và $y_0=0$, suy ra phương trình tiếp tuyến là $y=x-2$.
		Với $x_0=-2$, ta có $k=-7$ và $y_0=12$, suy ra phương trình tiếp tuyến là $y=-7x-2$.
		}
\end{vd}

\begin{dang}{Các bài toán vận dụng, thực tiễn}
\end{dang}

\begin{vd}%[TeX hóa SBT 11]%[1D2B2-6]
	Năm $2010$, dân số ở một tỉnh $D$ là $1\,038\,229$ người. Tính đến năm $2015$, dân số tỉnh đó là $1\,153\,600$ người. Cho biết dân số của tỉnh D được ước tính theo công thức $S(N)=A\mathrm{e}^{Nr}$ (trong đó $A$ là dân số của năm lấy làm mốc, $S$ là dân số sau $N$ năm, $r$ là tỉ lệ tăng dân số hàng năm được làm tròn đến hàng phần nghìn). Tốc độ gia tăng dân số (người/năm) vào thời điểm sau $N$ năm kể từ năm $2010$ được xác định bởi hàm số $S'(N)$. Tính tốc độ gia tăng dân số của tỉnh $D$ vào năm $2023$ (làm tròn kết quả đến hàng đơn vị theo đơn vị người/năm), biết tỉ lệ tăng dân số hàng năm không đổi.
	\loigiai{
		Tính từ năm $2010$ đến năm $2015$, chọn năm $2010$ làm mốc, ta có
		$$1\,153\,600=1\,038\,229\cdot\mathrm{e}^{5r}\Rightarrow r\approx 0{,}021.$$
		Khi đó, ta có $S(N)\approx 1\,038\,229\cdot \mathrm{e}^{0{,}021N}$, suy ra tốc độ gia tăng dân số vào thời điểm sau $N$ năm kể từ năm $2010$ là
		$$S'(N)\approx 0{,}021\cdot 1\,038\,229\cdot\mathrm{e}^{0{,}021N}\approx 21\,802{,}809\cdot \mathrm{e}^{0{,}021N}.$$
		Tốc độ gia tăng dân số tỉnh $D$ vào năm $2023$ (sau $13$ năm từ năm $2010$) là
		$$S'(13)\approx 21\,802{,}809\cdot \mathrm{e}^{0{,}021\cdot 13}\approx 28\,647 \text{ (người/năm)}.$$
	}
\end{vd} \dongcham{10}

\begin{vd}%[TeX hóa SBT 11]%[1D2B2-6]
	Một chất điểm chuyển động theo phương trình $s(t)=\dfrac{1}{3}t^3-3t^2+8x+2$, trong đó $t>0$, $t$ tính bằng giây, $s(t)$ tính bằng mét. Tính vận tốc tức thời của chất điểm tại thời điểm $t=5$ (s).
	\loigiai{
		Vận tốc tức thời của chất điểm tại thời điểm $t$ (s) là
		$$v(t)=s'(t)=t^2-6t+8.$$
		Vận tốc tức thời của chất điểm tại thời điểm $t=5$ (s) là
		$$v(5)=5^2-6\cdot 5+8=3\text{ (m/s)}.$$
	}
\end{vd} \dongcham{8}

\begin{vd}%[1D5B2-6]
	Nếu số lượng sản phẩm sản xuất được của một nhà máy là $x$ (đơn vị: trăm sản phẩm) thì lợi nhuận sinh ra là $P(x)=-200x^2+ 12800x-74000$ (nghìn đồng). Tính tốc độ thay đổi lợi nhuận của nhà máy đó khi sản xuất  $1200$ sản phẩm.
	\loigiai{
		Ta có $P'(x)=-2 \cdot 200 x + 12800 = -400x+12800$.\\
		Tốc độ thay đổi lợi nhuận của nhà máy đó khi sản xuất $1200$ sản phẩm là $P'(12)= -400 \cdot 12 + 12800 = 8000$.
	}
\end{vd} \dongcham{8}

\subsection{BÀI TẬP TỰ LUYỆN}

\begin{bt}
	Giải phương trình $f'(x)=0$, biết $f(x)$ được cho bởi công thức sau:
	\begin{tasks}(1)
		\task $f(x)=x^3-3x^2+2$;
		\task $f(x)=x^4-4x^2+2$;
		\task $f(x)=\dfrac{x^2+3x+3}{x+1}$;
		\task $f(x)=(x-1)\sqrt{2x+1}$;
		\task $f(x)=\sin x - \sqrt{3} \cos x$;
		\task $f(x)=\sin^2 x -x$.
	\end{tasks}
	\loigiai{
		\begin{enumerate}[\faCheckSquareO]
			\item Ta có $f'(x)=3x^2-6x$.\\
			Khi đó, $f'(x)=0\Leftrightarrow 3x^2-6x=0\Leftrightarrow\hoac{& x=0 \& x=2.}$\\
			Vậy phương trình $f'(x)=0$ có hai nghiệm $x=0$, $x=2$.
			\item $f'(x)=4x^3-8x = 4x\left(x^2-2 \right)$. \\
			$f'(x)=0 \Leftrightarrow x=0$ hoặc $x=\pm \sqrt{2}$.
			\item Ta có $f'(x)=\dfrac{x^2+2x}{(x+1)^2}$.\\
			Điều kiện $x\neq -1\quad (*)$.\\
			Khi đó, $f'(x)=0\Leftrightarrow x^2+2x=0\Leftrightarrow\hoac{& x=-2 \& x=0.}$\\
			Đối chiếu với điều kiện $(*)$, phương trình $f'(x)=0$ có hai nghiệm $x=-2$, $x=0$.
			\item 	Ta có $f'(x)=1\cdot\sqrt{2x+1}+(x-1)\cdot\dfrac{(2x+1)'}{2\sqrt{2x+1}}=\sqrt{2x+1}+\dfrac{x-1}{\sqrt{2x+1}}=\dfrac{2x+1+x-1}{\sqrt{2x+1}}=\dfrac{3x}{\sqrt{2x+1}}$.\\
			Điều kiện $x>-\dfrac{1}{2}\quad (*)$.\\
			Khi đó, $f'(x)=0\Leftrightarrow 3x=0\Leftrightarrow x=0$.\\
			Đối chiếu với điều kiện $(*)$, phương trình $f'(x)=0$ có một nghiệm $x=0$.
			\item $f'(x)=\cos x + \sqrt{3} \sin x =0 \Leftrightarrow \tan x =-\dfrac{1}{\sqrt{3}} \Leftrightarrow x= - \dfrac{\pi}{6}+k \pi$.
			\item $f'(x)=\sin 2x-1 =0 \Leftrightarrow \sin 2 x =1 \Leftrightarrow x=  \dfrac{\pi}{4}+k \pi$.
		\end{enumerate}
	}
\end{bt}


\begin{bt}%[TeX hóa SBT 11]%[1D5B2-3]
	Cho hàm số $y=x^2+3x$ có đồ thị $(C)$. Viết phương trình tiếp tuyến của đồ thị $(C)$ 
	\begin{tasks}(1)
		\task tại điểm có hoành độ bằng $-1$;
		\task tại điểm có tung độ bằng $4$.
		\task tiếp tuyến có hệ số góc bằng $1$.
		\task tiếp tuyến đi qua điểm $A(0;-5)$.
	\end{tasks} 
	\loigiai{
		Gọi $M(x_0;y_0)$ là tiếp điểm và $y'=2x+3$.
		\begin{enumerate}[a)]
			\item 
			Ta có $x_0=-1$ suy ra $y_0=-2$.\\
			Khi đó $y'(-1)=2\cdot (-1)+3=1$.\\
			Vậy phương trình tiếp tuyến là $$y=y'(-1)\cdot (x+1)-2=x-1.$$
			\item 
			Ta có $y_0=4\Leftrightarrow x_0^2+3x_0=4\Leftrightarrow \hoac{&x=1\\&x=-4.}$\\
			Với $x=1$ suy ra $y'(1)=5$.\\
			Phương trình tiếp tuyến là $$y=y'(1)\cdot (x-1)+4=5(x-1)+4=5x-1.$$
			Với $x=-4$ suy ra $y'(-4)=-5$.\\
			Phương trình tiếp tuyến là $$y=y'(-4)\cdot (x+4)+4=-5(x+4)+4=-5x-16.$$
		\end{enumerate}
	}
\end{bt}
\begin{bt}%[TeX hóa SBT 11]%[1D5B2-5]
	Cho hàm số $y=\dfrac{x-3}{x+2}$ có đồ thị $(C)$. Viết phương trình tiếp tuyến của đồ thị $(C)$ trong mỗi trường hợp sau
	\begin{tasks}
		\task $d$ song song với đường thẳng $y=5x-2$;
		\task $d$ vuông góc với đường thẳng $y=-20x+1$.
	\end{tasks}
	\loigiai{
		\begin{enumerate}[a)]
			\item 	Gọi $M(x_0;y_0)$ là tiếp điểm.\\
			Tiếp tuyến $d$ song song với đường thẳng $y=5x-2$ nên
			$$f'(x_0)=5\Leftrightarrow \dfrac{5}{(x_0+2)^2}=5\Leftrightarrow \hoac{&x_0=-1\\&x_0=-3.}$$
			Với $x_0=-1$ thì $y(-1)=-4$, phương trình tiếp tuyến $d$ là
			$$y=f'(x_0)\cdot (x-x_0)+y_0=5(x+1)-4=5x+1.$$
			Với $x_0=-3$ thì $y(-3)=6$, phương trình tiếp tuyến $d$ là 
			$$y=f'(x_0)\cdot (x-x_0)+y_0=5(x+3)+6=5x+21.$$
			\item Gọi $M(x_0;y_0)$ là tiếp điểm.\\
			Tiếp tuyến $d$ vuông góc với đường thẳng $y=-20x+1$ nên 
			$$f'(x_0)=\dfrac{1}{20}\Leftrightarrow \dfrac{5}{(x_0+2)^2}=\dfrac{1}{20}\Leftrightarrow \hoac{&x_0=-12\\&x_0=8.}$$
			Với $x_0=-12$ thì $y(-12)=\dfrac{3}{2}$, phương trình tiếp tuyến $d$ là
			$$y=f'(x_0)\cdot (x-x_0)+y_0=\dfrac{1}{20}(x+12)+\dfrac{3}{2}=\dfrac{1}{20}x+\dfrac{21}{10}.$$
			Với $x_0=8$ thì $y(8)=\dfrac{1}{2}$, phương trình tiếp tuyến $d$ là
			$$y=f'(x_0)\cdot (x-x_0)+y_0=\dfrac{1}{20}(x-8)+\dfrac{1}{2}=\dfrac{1}{20}x+\dfrac{1}{10}.$$
		\end{enumerate}
	}
\end{bt}


\begin{bt}%Câu 8
	Nếu số lượng sản phẩm sản xuất được của một nhà máy là $x$ (đơn vị: trăm sản phẩm) thì lợi nhuận sinh ra là $P(x)=200\left(x-2\right)\left(17-x\right)$ (nghìn đồng). Tính tốc độ thay đổi lợi nhuận của nhà máy đó khi sản xuất $3000$ sản phẩm.
	\loigiai{
		Ta có 
		$$P'(x)=200\left(x-2\right)'\left(17-x\right)+200\left(x-2\right){\left(17-x\right)'}=200\left(17-x\right)-200\left(x-2\right)=200\left(19-2x\right).$$
		Tốc độ thay đổi lợi nhuận của nhà máy đó khi sản xuất $3000$ sản phẩm là
		$$P'\left(30\right)=200\left(19-2\cdot 30\right)=-8200.$$}
\end{bt}

	\begin{bt}%[1D5K2-6]
	Một vật được phóng thẳng đứng lên trên từ mặt đất với vận tốc ban đầu là $v_0~(\mathrm{m} / \mathrm{s})$ (bỏ qua sức cản của không khí) thì độ cao $h$ của vật (tính bằng mét) sau $t$ giây được cho bởi công thức $h=v_0 t-\dfrac{1}{2} g t^2$ ($g$ là gia tốc trọng trường). Tìm vận tốc của vật khi chạm đất.
	\loigiai{
		Tại thời điểm vật chạm đất: $h=v_0 t-\dfrac{1}{2} g t^2=0(t>0)$.\\
		Giải phương trình ta được $t=\dfrac{2 v_0}{g}$.\\
		Vận tốc của vật khi chạm đất là $v=h'\left(\dfrac{2 v_0}{g}\right)=-v_0$.	
	}
\end{bt}

\begin{bt}%[1D5K2-6]
	Chuyển động của một hạt trên một dây rung được cho bởi công thức $s(t)=10+\sqrt{2} \sin \left(4 \pi t+\dfrac{\pi}{6}\right)$, trong đó $s$ tính bằng centimét và $t$ tính bằng giây. Tính vận tốc của hạt sau $t$ giây. Vận tốc cực đại của hạt là bao nhiêu? (Làm tròn kết quả đến chữ số thập phân thứ nhất).
	\loigiai{
		Vận tốc của hạt sau $t$ giây là: $v(t)=s'(t)=4 \pi \sqrt{2} \cos \left(4 \pi t+\dfrac{\pi}{6}\right)$.\\
		Vận tốc cực đại của hạt là: $v_{\max}=4 \pi \sqrt{2} \approx 17,8 \mathrm{~m} / \mathrm{s}$, đạt được khi: $\left|\cos \left(4 \pi t+\dfrac{\pi}{6}\right)\right|=1$ hay $t=\dfrac{5}{24}+\dfrac{k}{4}, k \in \mathbb{N}.$	
	}
\end{bt}

\begin{bt}%[TeX hóa SBT 11]%[1D2B2-6]
	Một mạch dao động điện từ $LC$ có lượng điện tích dịch chuyển qua tiết diện thẳng của một dây xác định bởi hàm số $Q(t)=10^{-5}\sin \left(2\,000t+\dfrac{\pi}{3}\right)$, trong đó $t>0$, $t$ tính bằng giây, $Q$ tính bằng Coulomb. Tính cường độ dòng điện tức thời $I$ (A) trong mạch tại thời điểm $t=\dfrac{\pi}{1\,500}$ (s). Biết $I(t)=Q'(t)$.
	\loigiai{
		Cường độ dòng điện tức thời trong mạch tại thời điểm $t$ (s) là
		$$I(t)=Q'(t)=10^{-5}\cdot 2\,000 \cos \left(2\,000t+\dfrac{\pi}{3}\right)=0{,}02\cos \left(2\,000t+\dfrac{\pi}{3}\right).$$
		Cường độ dòng điện tức thời trong mạch tại thời điểm $t=\dfrac{\pi}{1\,500}$ (s) là
		$$I\left(\dfrac{\pi}{1\,500}\right)=0{,}02\cos \left(2\,000\cdot \dfrac{\pi}{1\,500}+\dfrac{\pi}{3}\right)=0{,}01\text{ (A)}.$$
	}
\end{bt}

\begin{bt}%[TeX hóa SBT 11]%[1D2B2-6]
	Một tài xế đang lái xe ô tô, ngay khi phát hiện có vẫn cản phía trước đã phanh gấp lại nhưng vẫn xảy ra va chạm, chiếc ô tô để lại vết trượt dài $20{,}4$ m (được tính từ lúc bắt đầu đạp phanh đến khi xảy ra va chạm). Trong quá trình đạp phanh, ô tô chuyển động theo phương trình $s(t)=20t-\dfrac{5}{2}t^2$, trong đó $s(m)$ là độ dài quãng đường đi được sau khi phanh, $t$ (s) là thời gian tính từ lúc bắt đầu phanh ($0\leq t\leq 4)$.
	\begin{enumerate}[a)]
		\item Tính vận tốc tức thời của ô tô ngay khi đạp phanh. Hãy cho biết xe ô tô trên có chạy quá tốc độ hay không, biết tốc độ giới hạn cho phép là $70$ km/h.
		\item Tính vận tốc tức thời của ô tô ngay khi xảy ra va chạm?
	\end{enumerate}
	\loigiai{
		\begin{enumerate}[a)]
			\item Vận tốc tức thời của ô tô tại thời điểm $t$ (s) là $v(t)=s'(t)=20-5t$.\\
			Vận tốc tức thời của ô tô ngay khi đạp phanh ($t=0$ (s)) là
			$$v(0)=20-5\cdot 0=20\text{ (m/s)}.$$
			Ta có $20$ m/s$=72$ km/h $>70$ km/h. Suy ra ô tô trên đã chạy quá tốc độ giới hạn cho phép.
			\item Khi xảy ra va chạm, ô tô đã đi được $20{,}4$ m kể từ khi đạp phanh nên
			$$20{,}4=20t-\dfrac{5}{2}t^2\Leftrightarrow \hoac{&t=1{,}2\\&t=6{,}8.}$$
			Vì $0\le t\le 4$ nên $t=1{,}2$ (s).\\
			Vận tốc tức thời của ô tô ngay khi xảy ra va chạm ($t=1{,}2$ (s)) là
			$$v(1{,}2)=20-5\cdot 1{,}2=14\text{ (m/s)}.$$
		\end{enumerate}
	}
\end{bt}

